\documentclass[10pt,a4paper]{amsart}
\usepackage[margin=19mm]{geometry}
\usepackage{amsmath,amssymb,mathtools}
\usepackage[scr=rsfs,cal=euler]{mathalfa}
\usepackage{xfrac}
\usepackage[unicode,hidelinks,bookmarks=false]{hyperref}
\newcommand{\ie}{i.e.\ }
\newcommand{\cf}{cf.\ }
\newcommand{\resp}{respectively\ }
\newcommand{\Subsection}{\subsection}
\providecommand{\nobreakdash}{\nobreak\hbox{-}}
\setlength{\emergencystretch}{2em}
\multlinegap0pt
\usepackage{amssymb}
\usepackage{mathtools}
\newtheorem{lemma}{Lemma}
\newcommand{\Q}{\mathbb{Q}}
\title{Edge-Span Chern Algebras of Graphical Configuration Spaces}
\author{Jackson Walters}
\date{Source: arXiv:2607.29483v1 (CC BY 4.0)}
\begin{document}
\maketitle

\noindent\emph{Source and licence.} Edge-Span Chern Algebras of Graphical Configuration Spaces. Version arXiv:2607.29483v1 (CC BY 4.0), distributed by arXiv under the Creative Commons Attribution 4.0 International licence, http://creativecommons.org/licenses/by/4.0/, as recorded in the arXiv licence metadata for this exact version and shown on the abstract page https://arxiv.org/abs/2607.29483v1. Full licence notice: \texttt{licenses/arxiv-2607.29483-CC-BY-4.0.txt}.\par\smallskip

\noindent\textit{Editorial scope.} Counted unit: the article's lemma lem:bipartite-conductor (source0.tex lines 1077--1093). The article's complete original proof of it is delivered with it (source0.tex lines 1095--1142, one \texttt{proof} environment of 48 lines). No mathematics is written, repaired or re-derived: the statement and the proof are the article's own lines, in the article's own notation, and no retained line is substituted or re-flowed.  No further source-local context is needed: every cross-reference of every retained line resolves inside the two delivered spans.  Nothing was omitted from any retained span, so the package carries no omission record.  No retained line contains a citation, so the wrapper carries no bibliography and no bibliography entry.  No retained line contains a figure environment, an image-inclusion command or drawing code, so no image is delivered and the package declares no asset dependency.  Editorial formatting only, in addition: an AMS \texttt{amsart} preamble that restates the article's own declarations and macros used by the retained lines, namely the environments \texttt{lemma} on separate counters, together with the standard packages those lines need.  The retained environments print in this document as Lemma 1 in order of appearance, whereas the article prints the same items with the numbering of its own full document.  The complete native source of the article as distributed in the arXiv e-print for this exact version is bundled unchanged as \texttt{sources/206490/source0.tex}.
\begin{lemma}[Bipartite conductor lemma]\label{lem:bipartite-conductor}
Let $G$ be a connected bipartite graph with at least one edge and color
classes $V_0,V_1$.  Define
\[
  \pi_G:P_G\longrightarrow\Q[x,y],
  \qquad
  h_v\longmapsto
  \begin{cases}
    x,&v\in V_0,\\
    y,&v\in V_1.
  \end{cases}
\]
Then
\begin{equation}\label{eq:polynomial-bipartite-preimage}
  S_G=\pi_G^{-1}\bigl(\Q[x+y,xy]\bigr).
\end{equation}
\end{lemma}

\begin{proof}
Fix a vertex $u$.  The edge sums along paths express every $h_v$ in terms of
$S_G$ and $h_u$, so $P_G=S_G[h_u]$.  If $uw$ is any edge incident to $u$,
then
\[
  h_u^2-(h_u+h_w)h_u+h_uh_w=0,
\]
and multiplying this monic relation by $h_u^{k-2}$ gives, for $k\geq2$,
\[
  h_u^k=(h_u+h_w)h_u^{k-1}-h_uh_w\,h_u^{k-2}.
\]
Both coefficients on the right belong to $S_G$, so induction on $k$ reduces
every power of $h_u$ to an $S_G$-linear combination of $1$ and $h_u$.
Together with $P_G=S_G[h_u]$, this gives
\begin{equation}\label{eq:polynomial-rank-two}
  P_G=S_G+S_Gh_u.
\end{equation}

Consider a path of length two, $u-w-v$, and set $\delta=h_u-h_v$.  The edge
generators give
\[
  \delta=(h_u+h_w)-(h_w+h_v)\in S_G,
\]
\[
  h_w\delta=h_uh_w-h_wh_v\in S_G,
  \qquad
  h_u\delta=(h_u+h_w)\delta-h_w\delta\in S_G.
\]
Equation \eqref{eq:polynomial-rank-two} therefore implies
$\delta P_G\subseteq S_G$.  In other words, every difference arising from a
length-two path belongs to the conductor
\[
  \mathfrak c_G=\{f\in P_G:fP_G\subseteq S_G\}.
\]
Any two vertices in the same color class are joined by an even path, and
their difference is a sum of such length-two differences.  Consequently,
\[
  J_G=\ker\pi_G
  =(h_u-h_v:u,v\in V_0\text{ or }u,v\in V_1)
  \subseteq\mathfrak c_G\subseteq S_G.
\]

Certainly $\pi_G(S_G)=\Q[x+y,xy]$.  Conversely, if
$\pi_G(f)=F(x+y,xy)$, choose one edge $ab$ and set
$g=F(h_a+h_b,h_ah_b)$.  Then $g\in S_G$ and
$f-g\in J_G\subseteq S_G$, proving
\eqref{eq:polynomial-bipartite-preimage}.
\end{proof}

\end{document}
